% ============================================================
%  Preámbulo — plantilla main-udea (tema default + azul)
%  adaptada al contenido de los Lemas B y C (grupo fundamental)
% ============================================================

% ---------- Idioma y codificación ----------
\usepackage[T1]{fontenc}
\usepackage[utf8]{inputenc}
\usepackage[spanish]{babel}

% ---------- Matemáticas ----------
\usepackage{amsmath}
\usepackage{amssymb}
\usepackage{amsfonts}   % \varnothing
\usepackage{amsthm}
\usepackage{mathtools}

% ---------- Tablas, gráficos y utilidades ----------
\usepackage{array}
\usepackage{booktabs}
\usepackage{graphicx}
\usepackage{ragged2e}   % \justifying
\usepackage{hyperref}
\usepackage{minted}
\usepackage{csquotes}   % must come after fvextra (loaded by minted)

% Suppress \\ in PDF author/title strings
\pdfstringdefDisableCommands{\let\\=\space}

% ---------- TikZ ----------
\usepackage{tikz}
\usetikzlibrary{arrows.meta, decorations.pathreplacing, patterns, calc, positioning}

% ---------- Bibliografía ----------
\usepackage[backend=biber, style=numeric]{biblatex}
\addbibresource{references.bib}

% ============================================================
%  Tema y colores (plantilla main-udea)
% ============================================================

% ---------- Paleta de colores ----------
% Los dos morados estructurales se re-mapean al azul de la plantilla;
% los acentos (verde/coral/ámbar) se mantienen para armonizar con el azul.
\definecolor{morado}{HTML}{1A73E8}        % azul de plantilla (antes púrpura)
\definecolor{moradoclaro}{HTML}{CCE5FF}   % azul claro
\definecolor{verdeTeal}{RGB}{15,110,86}
\definecolor{verdeClaro}{RGB}{225,245,238}
\definecolor{coralOscuro}{RGB}{153,60,29}
\definecolor{ambar}{RGB}{133,79,11}
\definecolor{grisClaro}{RGB}{245,244,240}

\usetheme{default}                                        % tema simple y limpio
\setbeamercolor{frametitle}{fg=morado}                    % títulos de slide
\setbeamercolor{structure}{fg=morado}                     % viñetas, toc, etc.
\setbeamercolor{title in head/foot}{fg=morado}
\setbeamercolor{date in head/foot}{fg=black}

% Pie de página con número de slide (plantilla main-udea)
\setbeamertemplate{navigation symbols}{}
\addtobeamertemplate{navigation symbols}{}{%
    \usebeamerfont{footline}%
    \usebeamercolor[fg]{footline}%
    \hspace{1em}%
    \insertframenumber/\inserttotalframenumber
}

% ---------- Colores de bloques ----------
\setbeamercolor{block title}{bg=morado, fg=white}
\setbeamercolor{block body}{bg=grisClaro, fg=black}
\setbeamercolor{block title alerted}{bg=coralOscuro, fg=white}
\setbeamercolor{block body alerted}{bg=grisClaro, fg=black}
\setbeamercolor{block title example}{bg=verdeTeal, fg=white}
\setbeamercolor{block body example}{bg=verdeClaro, fg=black}
\setbeamercolor{section in toc}{fg=morado}

% ---------- Comandos matemáticos ----------
\newcommand{\R}{\mathbb{R}}
\newcommand{\Z}{\mathbb{Z}}
\newcommand{\I}{[0,1]}
\newcommand{\IxI}{I \times I}
\newcommand{\pione}{\pi_1}
\newcommand{\xo}{x_0}
\newcommand{\simxo}{\sim_{\xo}}
\newcommand{\abar}{\bar{\alpha}}
\newcommand{\clase}[1]{[#1]}

% ---------- Separadores de sección ----------
\AtBeginSection[]{%
  \begin{frame}{Contenido}
    \tableofcontents[currentsection]
  \end{frame}%
}
